\documentclass[a4paper,10pt]{article}
\usepackage[czech]{babel}
\usepackage[T1]{fontenc}
\usepackage[utf8]{inputenc}
\usepackage{amsmath,amssymb}
\usepackage[margin=18mm]{geometry}
\setlength{\parindent}{0pt}

\begin{document}
\section*{Komplexní čísla C/D -- řešení}
Podklad pro učitele k testu \texttt{02\_test\_KomplexniCislaZaverecny\_CD.tex}, 3.~ročník SPŠ.

Všechny matematické údaje pocházejí ze schválených úloh učebnice:\\
\texttt{ucebnice/06\_Komplexni\_cisla/Komplexní čísla.tex}.
Výběr a přiřazení k variantám schválil učitel v této konverzaci.
Ve čtvrté úloze obou variant bylo na jeho pokyn číslo $z_1$ převedeno
do algebraického tvaru; hodnoty čísel i požadované výpočty zůstaly zachovány.

\begin{center}
\begin{tabular}{c|ll}
Úloha testu & Varianta C: zdroj & Varianta D: zdroj\\ \hline
1 & Úloha 1.1 c), řádek 4525 & Úloha 1.1 i), řádek 4531\\
2 & Úloha 2.22 h), řádek 4927 & Úloha 2.22 i), řádek 4928\\
3 & Úloha 4.49 b), řádek 5451 & Úloha 4.49 c), řádek 5452\\
4 & Úloha 3.35 b), řádek 5162 & Úloha 3.35 d), řádek 5166\\
5 & Úloha 3.44 a), řádek 5325 & Úloha 3.44 d), řádek 5346
\end{tabular}
\end{center}
Řádky odpovídají zdroji při sestavení; rozhodující je číslo úlohy a písmeno.
Zařazení odpovídá tematickému plánu
\texttt{3\_rocnik/tematicke\_plany/SPS.pdf}, s.~1, body 1.1--1.5.

\subsection*{Varianta C}
\begin{enumerate}
\item Vypočítáme součin postupným násobením:
\[
(2+i)(3+i)(1+2i)=(5+5i)(1+2i)=5+15i-10=-5+15i.
\]
Kontrola jiným seskupením:
$(3+i)(1+2i)=1+7i$, tedy $(2+i)(1+7i)=-5+15i$.

\item Protože $1-2i\neq0$ a $|(2+i)^2|=|2+i|^2$, dostáváme
\[
\frac{|(2+i)^2|}{|1-2i|}
=\frac{(\sqrt5)^2}{\sqrt5}=\sqrt5.
\]
Kontrola přes algebraický tvar podílu:
\[
\frac{(2+i)^2}{1-2i}=\frac{(3+4i)(1+2i)}5=-1+2i,
\qquad |-1+2i|=\sqrt5.
\]

\item Pro rovnici $5x^2-2x+1=0$ je diskriminant $D=4-20=-16$.
\[
x_{1,2}=\frac{2\pm4i}{10}=\frac{1\pm2i}{5},
\qquad K=\left\{\frac{1-2i}{5};\frac{1+2i}{5}\right\}.
\]
Zkouška pro oba kořeny:
\[
5\left(\frac{1\pm2i}{5}\right)^2
-2\frac{1\pm2i}{5}+1
=\frac{-3\pm4i-2\mp4i+5}{5}=0.
\]
Jsou to dva různé kořeny kvadratické rovnice, tedy všechny její kořeny v $\mathbb C$.

\item Je dáno $z_1=1-\sqrt3i$ a
$z_2=4(\cos\frac{\pi}{6}+i\sin\frac{\pi}{6})$.
Číslo $z_1$ má modul $\sqrt{1+3}=2$ a leží ve IV.~kvadrantu,
tedy $z_1=2(\cos\frac{5\pi}{3}+i\sin\frac{5\pi}{3})$. Počítáme:
\begin{align*}
z_1z_2
&=8\left(\cos\frac{11\pi}{6}+i\sin\frac{11\pi}{6}\right)
=4\sqrt3-4i,\\
\frac{z_1}{z_2}
&=\frac12\left(\cos\frac{3\pi}{2}+i\sin\frac{3\pi}{2}\right)
=-\frac12i.
\end{align*}
U součinu argumenty sčítáme, u podílu je odečítáme;
$\frac{5\pi}{3}+\frac\pi6=\frac{11\pi}{6}$,
$\frac{5\pi}{3}-\frac\pi6=\frac{3\pi}{2}$.
Kontrola v algebraickém tvaru: $z_1=1-\sqrt3i$, $z_2=2\sqrt3+2i$.
Platí $(1-\sqrt3i)(2\sqrt3+2i)=4\sqrt3-4i$
a $(-\frac12i)(2\sqrt3+2i)=1-\sqrt3i=z_1$.

\item Pro $z=\frac12(\cos\frac{4\pi}{9}+i\sin\frac{4\pi}{9})\neq0$
nejprve použijeme Moivreovu větu a poté převrátíme výsledek:
\[
z^3=\frac18\left(\cos\frac{4\pi}{3}+i\sin\frac{4\pi}{3}\right),
\qquad
\frac1{z^3}=8\left(\cos\frac{2\pi}{3}+i\sin\frac{2\pi}{3}\right)
=-4+4\sqrt3i.
\]
Při převrácení se modul převrátí a argument změní znaménko;
$-\frac{4\pi}{3}+2\pi=\frac{2\pi}{3}$.
Kontrola: $z^3=\frac{-1-\sqrt3i}{16}$ a
$\frac{-1-\sqrt3i}{16}(-4+4\sqrt3i)=1$.
\end{enumerate}

\newpage
\section*{Komplexní čísla -- varianta D}
Řešení k testu \texttt{02\_test\_KomplexniCislaZaverecny\_CD.tex}.
\begin{enumerate}
\item Nejprve vypočítáme druhou mocninu a poté násobíme:
\[
(1-2i)^3=(-3-4i)(1-2i)=-11+2i.
\]
Kontrola vzorcem pro třetí mocninu:
$1-6i+12i^2-8i^3=1-6i-12+8i=-11+2i$.

\item Protože $2+i\neq0$, použijeme vlastnosti absolutní hodnoty:
\[
\left|\frac{(1-i)^3}{2+i}\right|
=\frac{|1-i|^3}{|2+i|}
=\frac{(\sqrt2)^3}{\sqrt5}=\frac{2\sqrt{10}}5.
\]
Kontrola přes algebraický tvar podílu:
\[
\frac{(1-i)^3}{2+i}=\frac{(-2-2i)(2-i)}5=-\frac65-\frac25i,
\qquad
\sqrt{\frac{36}{25}+\frac4{25}}=\frac{2\sqrt{10}}5.
\]

\item Pro rovnici $3x^2-2x+1=0$ je diskriminant $D=4-12=-8$.
\[
x_{1,2}=\frac{2\pm2\sqrt2i}{6}=\frac{1\pm\sqrt2i}{3},
\qquad K=\left\{\frac{1-\sqrt2i}{3};\frac{1+\sqrt2i}{3}\right\}.
\]
Zkouška pro oba kořeny:
\[
3\left(\frac{1\pm\sqrt2i}{3}\right)^2
-2\frac{1\pm\sqrt2i}{3}+1
=\frac{-1\pm2\sqrt2i-2\mp2\sqrt2i+3}{3}=0.
\]
Jsou to dva různé kořeny kvadratické rovnice, tedy všechny její kořeny v $\mathbb C$.

\item Je dáno $z_1=-\frac{5\sqrt3}{2}-\frac52i$ a
$z_2=\cos\frac\pi3+i\sin\frac\pi3$.
Číslo $z_1$ má modul $\sqrt{\frac{75}{4}+\frac{25}{4}}=5$ a leží ve III.~kvadrantu,
tedy $z_1=5(\cos\frac{7\pi}{6}+i\sin\frac{7\pi}{6})$. Počítáme:
\begin{align*}
z_1z_2
&=5\left(\cos\frac{3\pi}{2}+i\sin\frac{3\pi}{2}\right)=-5i,\\
\frac{z_1}{z_2}
&=5\left(\cos\frac{5\pi}{6}+i\sin\frac{5\pi}{6}\right)
=-\frac{5\sqrt3}{2}+\frac52i.
\end{align*}
U součinu používáme součet $\frac{7\pi}{6}+\frac\pi3=\frac{3\pi}{2}$,
u podílu rozdíl $\frac{7\pi}{6}-\frac\pi3=\frac{5\pi}{6}$.
Kontrola v algebraickém tvaru:
$z_1=-\frac{5\sqrt3}{2}-\frac52i$, $z_2=\frac12+\frac{\sqrt3}{2}i$.
Přímé násobení dává $z_1z_2=-5i$ a
\[
\left(-\frac{5\sqrt3}{2}+\frac52i\right)
\left(\frac12+\frac{\sqrt3}{2}i\right)
=-\frac{5\sqrt3}{2}-\frac52i=z_1.
\]

\item Pro $z=\frac13(\cos\frac{5\pi}{6}+i\sin\frac{5\pi}{6})\neq0$
nejprve použijeme Moivreovu větu a poté převrátíme výsledek:
\[
z^2=\frac19\left(\cos\frac{5\pi}{3}+i\sin\frac{5\pi}{3}\right),
\qquad
\frac1{z^2}=9\left(\cos\frac\pi3+i\sin\frac\pi3\right)
=\frac92+\frac{9\sqrt3}{2}i.
\]
Argument převrácené hodnoty je
$-\frac{5\pi}{3}+2\pi=\frac\pi3$.
Kontrola: $z^2=\frac{1-\sqrt3i}{18}$ a
$\frac{1-\sqrt3i}{18}(\frac92+\frac{9\sqrt3}{2}i)=1$.
\end{enumerate}

\subsection*{Kontrola a srovnatelnost variant}
Všechny jmenovatele jsou nenulové, moduly kladné a oba kvadratické kořeny
v každé variantě byly zkontrolovány dosazením.
Součiny, podíly a převrácené mocniny byly ověřeny také přesnou aritmetikou;
goniometrická zadání nezávisle numericky z původních úhlů.
Všechny převzaté matematické výrazy odpovídají schválenému výběru.

Varianty mají shodné okruhy, počet úloh a přibližně srovnatelnou náročnost.
V D3 je navíc odmocnina v imaginární části kořenů, zatímco C1 obsahuje
součin tří různých dvojčlenů. Náročnost proti A/B zvyšují vícefázové
operace, mocniny uvnitř absolutní hodnoty, nenormované kvadratické rovnice,
součin i podíl ve čtvrté úloze a převrácené mocniny v páté úloze.
Orientační časy jsou 3--4, 3--5, 4--6, 6--8 a 4--6 minut;
celkem přibližně 25--35 minut včetně kontroly. Odhad nebyl ověřen ve třídě.
\end{document}
